\section{Formal certification: progressive thirteen-stage proof ledger}\label{sec:formal}
A pinned Lean~4/mathlib~v4.24.0 project certifies proof-critical
identities of Lemma~\ref{lem:gram}; sources and lockfile reside in the
\href{https://github.com/tonynewton79-web/newton-400-conjectures}{public
Newton Astro Labs repository}. The full graph-spectral theorem is
\emph{not yet kernel-certified}.

\paragraph{Stage 1: rational margins and Gram positivity.}
\texttt{NewtonConjectures.lean} certifies the strict base-case margins
$506/11025$ and $1733/58800$, rational reciprocal second differences,
positive weights, and nonnegative weighted Gram quadratic forms.

\paragraph{Stage 2: universal reciprocal-distance kernel.}
\texttt{CycleKernel.lean} proves telescoping for all
$f:\mathbb N\to\mathbb Q$. For $1\le d<m$, the reciprocal
specialization is
\begin{align}\label{eq:lean-kernel}
\frac1d={}&\frac1m+
\sum_{j=d+1}^{m-1}(j-d)\frac{2}{(j-1)j(j+1)}\notag\\
&+\frac{m-d}{m(m-1)}.
\end{align}

\paragraph{Stage 3: exact cyclic Gram entries.}
\texttt{CyclicIntersection.lean}, \texttt{CyclicZMod.lean}
and \texttt{CyclicGram.lean} formalize genuine cyclic incidence
matrices over $\mathbb Z/n\mathbb Z$ and their \texttt{Matrix.mul}.
For all vertices $u,v$ and $2j\le n$, including wraparound,
disjoint arcs, and antipodes, the kernel proves
\begin{equation}\label{eq:lean-overlap}
(B_jB_j^{\mathsf T})_{uv}=\max\{j-d(u,v),0\}.
\end{equation}

\paragraph{Stage 4: actual weighted Gram diagonal.}
For every $m\ge2$, \texttt{GramDiagonal.lean} certifies
\begin{equation}\label{eq:lean-diag}
\frac1m+\sum_{j=2}^{m-1}\frac{2}{(j-1)(j+1)}
+\frac1{m-1}=\frac32.
\end{equation}
It also proves the identical diagonal for the \emph{actual weighted
incidence matrices} of every $C_n$, $n\ge4$.

\paragraph{Stage 5: all off-diagonal distances.}
For every $m\ge2$ and $1\le d\le m$,
\texttt{OffDiagonal.lean} proves
\begin{equation}\label{eq:lean-offdiag}
\frac1m+\sum_{j=2}^{m}w_j\max\{j-d,0\}=\frac1d.
\end{equation}
These are exactly the weights of \eqref{eq:weights}, with the
terminal case $d=m$ included using natural truncated subtraction.


\paragraph{Stage 6: the complete rational matrix equality.}
The new kernel-checked module \texttt{FullGram.lean} formally defines
the shifted reciprocal-distance matrix over $\mathbb Z/n\mathbb Z$,
with diagonal $3/2$ and off-diagonal entries $1/d(u,v)$.
For every integer $n\ge4$, the theorem
\texttt{cyclic\_shifted\_reciprocal\_eq\_gram} establishes
the \emph{actual equality of rational matrices}
\begin{equation}\label{eq:lean-fullgram}
R_n+\tfrac32I=\tfrac1mJ+\sum_{j=2}^m w_jB_jB_j^{\mathsf T},
\quad m=\lfloor n/2\rfloor.
\end{equation}
This result incorporates both parity branches, wraparound, all
off-diagonal distances and the diagonal normalization.
A separate Lean proof shows $1\le d(u,v)\le m$ whenever $u\ne v$.
The \href{https://github.com/tonynewton79-web/newton-400-conjectures/actions/runs/39133219628}{Stage~6 build}
completed 7,359 jobs, with 113 declarations passing the
trusted-foundation axiom audit.


\paragraph{Stage 7: kernel-certified matrix positivity.}
The module \texttt{SpectralPSD.lean} uses the exact Stage~6
matrix theorem and proves, for every $n\ge4$, that the
\emph{rational} matrix $R_n+\tfrac32I$ is positive semidefinite
in the formal \texttt{Matrix.PosSemidef} sense. The kernel
checks nonnegative piecewise weights, the constant rank-one
component, each actual incidence Gram product, and their
finite weighted sum. Thus
\begin{equation}\label{eq:lean-psd}
x^{\mathsf T}(R_n+\tfrac32I)x\ge0
\qquad(x\in\mathbb Q^n).
\end{equation}
The \href{https://github.com/tonynewton79-web/newton-400-conjectures/actions/runs/39134282319}{Stage~7 build}
passed 7,360 Lean jobs and audited 125 declarations within
the standard foundational axiom allowlist. The conversion to
a real-matrix/eigenvalue statement is a separate formal gate;
the analytic consequence for real symmetric matrices remains
available by the usual density and spectral arguments.


\paragraph{Stage 8: real Gram identity and spectral positivity.}
Two additional formal modules, \texttt{RealGram.lean} and
\texttt{RealSpectral.lean}, transport the actual weighted
cyclic-incidence Gram construction from $\mathbb Q$ to
$\mathbb R$. The kernel verifies both the equality with the
real shifted reciprocal-distance matrix and real positive
semidefiniteness. It then proves
\begin{equation}\label{eq:lean-spectrum}
\operatorname{spec}_{\mathbb R}(R_n+\tfrac32 I)
\subseteq[0,\infty)\qquad(n\ge4).
\end{equation}
Consequently the familiar lower bound
$\lambda_{\min}(R_n)\ge-3/2$ follows via the standard scalar
spectral shift; that last unshifted-eigenvalue formulation has
not yet been compiled as its own Lean theorem.
\href{https://github.com/tonynewton79-web/newton-400-conjectures/actions/runs/39139083900}{Stage~8 GitHub Actions}
passed 7,362 build jobs; 142 declarations survived the axiom
audit with only the standard foundational allowlist.


\paragraph{Stage 9: parity-uniform rational comparisons.}
In \texttt{HarmonicComparison.lean}, Lean certifies exact
harmonic-number recurrences and proves both rational spectral-spread
upper bounds decrease strictly with the cycle half-length.
Writing $q_o(m)=(4H_m+3)/m$ and
$q_e(m)=4(2H_{m-1}+1/m+3/2)/(2m-1)$, the kernel proves
\begin{align}\label{eq:lean-parity-bounds}
q_o(m)&<2-(22/105)^2 &&(m\ge7),\notag\\
q_e(m)&<2-(11/28)^2 &&(m\ge8).
\end{align}
The boundary evaluations $q_o(7)=468/245$ and
$q_e(8)=1907/1050$ are exact and are joined to the
universal monotonicity results, not inferred from finite samples.
The \href{https://github.com/tonynewton79-web/newton-400-conjectures/actions/runs/39140710038}{Stage~9 kernel run}
passed 7,363 build jobs and audited 180 declarations.
The precise trigonometric/adjacency eigenvalue comparison and finite
$n=3,\ldots,14$ complementary cases still require Lean treatment.


\paragraph{Stage 10: trigonometric and all-cycle closure.}
The modules \texttt{TrigonometricEndgame.lean} and
\texttt{AllCycleAnalytic.lean} use the mathlib inequality
$\cos x\ge1-x^2/2$ and $\pi<22/7$ to verify, without numerical
approximation,
\begin{align}\label{eq:lean-cos-comparison}
q_o(m)&<2\cos\!\frac{\pi}{2m+1} &&(m\ge7),\notag\\
q_e(m)&<2\cos\!\frac{\pi}{m} &&(m\ge8).
\end{align}
A further single Lean theorem selects the appropriate parity
and proves the closed-form strict comparison for every $n\ge15$.
\href{https://github.com/tonynewton79-web/newton-400-conjectures/actions/runs/39146022640}{The accepted Stage~10 run}
completed 7,365 build jobs and audited 197 declarations.
The actual cycle adjacency matrix over $\mathbb Z/n\mathbb Z$
is now defined in \texttt{AdjacencyDefinition.lean}, with its
symmetry and entry nonnegativity certified in the
\href{https://github.com/tonynewton79-web/newton-400-conjectures/actions/runs/39146480928}{204-declaration kernel audit}.
This isolates the remaining ordered-eigenvalue theorem.
This formally closes the \emph{analytical} comparison, but it
does not identify the right-hand cosine with the \emph{ordered}
second-smallest adjacency eigenvalue inside Lean.


\paragraph{Stage 11: final graph statement and exact bridge.}
The module \texttt{FinalTheoremProbe.lean} now defines the
\emph{actual} cycle gravity matrix by the published
$4/((n-1)d(u,v))$ entries and its numerical range as the
supremum minus infimum of the real spectrum. It also uses
mathlib's \emph{ordered} eigenvalues (including multiplicity)
of the actual cycle adjacency matrix, selecting the
second-smallest at index $n-2$. The companion module
\texttt{ActualGravityBridge.lean} kernel-proves
\begin{equation}\label{eq:actual-shift}
R_n+\tfrac32 I=\texttt{realShiftedReciprocalCycle}(n),
\end{equation}
with $R_n$ now defined by its zero-diagonal graph entries,
and certifies the positive semidefiniteness of this
\emph{actual} shifted graph matrix. It also verifies the
gravity scale factor $4/(n-1)$.

The \href{https://github.com/tonynewton79-web/newton-400-conjectures/actions/runs/39151152516}{Stage~11 Lean run}
passed 7,368 build jobs with 223 declarations admitted by the
standard foundational axiom audit. The end-to-end numerical
Graffiti inequality is accepted by the kernel \emph{only
conditionally}, assuming (i) identification of the actual
second-smallest cycle adjacency eigenvalue with
\eqref{eq:adj}, and (ii) the actual gravity spectral-span
bound \eqref{eq:spread}. The theorem's statement therefore
exposes rather than closes both missing spectral bridges.
The finite converse $C_3,\ldots,C_{14}$ is also pending
kernel certification. There is no claim of a complete
unconditional Lean proof.


\paragraph{Stage 12: distance-shell upper matrix certificate.}
The modules \texttt{DistanceShellLaplacian.lean},
\texttt{DistanceShellAdjacency.lean}, and
\texttt{DistanceShellDegreeBridge.lean} construct an
independent positive-semidefinite \emph{upper} matrix
certificate, complementary to the shifted-Gram lower
certificate. For each distance $1\le d\le m$, let
$E_d$ be the undirected distance-$d$ shell graph of $C_n$.
Writing $L(E_d)$ for its ordinary graph Laplacian and
$D_n$ for the weighted diagonal shell-degree matrix,
the kernel proves the exact identities
\begin{align}\label{eq:lean-laplacian}
R_n&=\sum_{d=1}^m\frac1d\,A(E_d),\notag\\
L_n&=\sum_{d=1}^m\frac1d\,L(E_d)=D_n-R_n.
\end{align}
Each graph Laplacian is positive semidefinite, so
\begin{equation}\label{eq:lean-gravity-upper}
D_n-R_n\succeq0,\qquad
\frac4{n-1}D_n-\Gr_n\succeq0\quad(n\ge4).
\end{equation}
These are \emph{actual} cycle-distance matrices and
the original gravity scaling, with no spectral target
substituted into the definition.
\href{https://github.com/tonynewton79-web/newton-400-conjectures/actions/runs/39153984478}{The Stage~12 kernel audit}
accepted 250 project declarations after 7,371 successful
build jobs. The remaining graph-specific step is to
prove inside Lean that $D_n=s_n I$, with
$s_n=2H_m$ for odd $n$ and
$s_n=2H_{m-1}+1/m$ for even $n$. Until that identity
and the spectral extraction compile, \eqref{eq:lean-gravity-upper}
is not by itself the full numerical-span inequality.


\paragraph{Stage 13: rotation-invariant shell counts and scalar degree.}
The new modules \texttt{DistanceShellDegreeRotation.lean} and
\texttt{DistanceShellScalarDegree.lean} give a graph-native
proof that distance-shell adjacency is preserved by simultaneous
cyclic translation of both endpoints. A finite bijection then
proves that each shell has a vertex-independent neighbor count.
The inverse-distance-weighted count $s_n$, evaluated at
the zero vertex, is therefore constant over the cycle, and
the kernel establishes the exact scalar-matrix identity
\begin{equation}\label{eq:lean-scalar-degree}
D_n=s_nI,\qquad
s_n:=\sum_{d=1}^{\lfloor n/2\rfloor}
\frac{N_d(0)}d.
\end{equation}
Together with the Stage~12 weighted Laplacian theorem,
this gives the \emph{actual} graph matrix inequality
\begin{equation}\label{eq:lean-scalar-upper}
s_nI-R_n\succeq0 \qquad(n\ge4).
\end{equation}
The \href{https://github.com/tonynewton79-web/newton-400-conjectures/actions/runs/39157366365}{successful Stage~13 Lean run}
compiled 7,373 build jobs and audited 268 project
declarations within the foundational axiom allowlist.
No numerical fitting or unproved regularity assumption enters
this scalar identity.

For odd $n=2m+1$, each nonzero distance $1\le d\le m$
has the two endpoints $u\pm d$, giving $s_n=2H_m$.
For even $n=2m$, the antipodal shell $d=m$ has one
endpoint while every smaller positive shell has two, giving
$s_n=2H_{m-1}+1/m$. These elementary arithmetic
identities were independently replayed with exact rational
arithmetic for all $4\le n\le160$ (685,875 shell-cardinality
checks and 12,874 vertex row-sum checks), and
\emph{have not yet been kernel-certified uniformly in $n$}.
Nor has the resulting full gravity numerical spectral-span
bound been extracted in Lean. The graph-theoretic
eigenvalue ordering and the $C_3,\ldots,C_{14}$ converse
also remain formal obligations.

\paragraph{Originality screen and claim calibration.}
The primary Roucairol--Cazenave source already refutes Graffiti~197
using $C_{17}$ and conjectures only the $C_{4r+1}$ family;
the present claimed advance is the stronger sharp classification
and constructive matrix certificate, not the first refutation.
Earlier work also studies reciprocal-distance (Harary) spectra
and their spreads. The focused source screen found no identical
all-$n$ threshold proof, but no search result establishes priority
or substitutes for external specialist review.
See the repository's
\href{https://github.com/tonynewton79-web/newton-400-conjectures/blob/main/cases/001/PRIOR_ART_AUDIT.md}{prior-art audit}.

\paragraph{Independent verification.}
\href{https://github.com/tonynewton79-web/newton-400-conjectures/actions/runs/39131413230}{GitHub Actions run 39131413230} passed 7,358 Lean build jobs.
Its axiom audit accepted 100 declarations with only
\texttt{propext}, \texttt{Classical.choice}, \texttt{Quot.sound};
the proof-placeholder check passed. Exact
\texttt{fractions.Fraction} arithmetic independently confirmed
all 7,259 cases $2\le m\le120$, $1\le d\le m$.
See the versioned \href{https://github.com/tonynewton79-web/newton-400-conjectures/blob/main/cases/001/LEAN_STATUS.md}{proof ledger}.

\paragraph{Claim boundary.}
Stages 1--12 now kernel-certify the exact shifted reciprocal
Gram decomposition and its real spectral positivity, the
odd/even harmonic and trigonometric analytic comparisons,
and the actual gravity upper Loewner bound
\eqref{eq:lean-gravity-upper}, and the scalar degree identity
\eqref{eq:lean-scalar-degree}. The strict Graffiti
inequality has been compiled only \emph{conditionally}
upon its actual spectral interfaces. The outstanding Lean
proofs are: the explicit parity-dependent harmonic evaluation of $s_n$,
gravity spectral-span extraction, the adjacency matrix's
ordered second-smallest eigenvalue and multiplicities, and
the finite nonviolating cases $C_3,\ldots,C_{14}$.
The full sharp-cycle theorem therefore remains a candidate
for end-to-end kernel certification and independent expert review.
